\chapter{The daemon}

\label{ch:daemon}

This chapter answers what the daemon is, what it refuses to start without, and where the real
security boundary sits. The short answer: one process per state directory, one writer per
store, a socket for everything else, and a filesystem confinement applied before the first
byte is served. The interesting part is not the wiring diagram; it is the order of the checks,
because each check exists to make a specific failure impossible rather than unlikely.

\section{The production shape}

Four binaries ship: \texttt{pii}, \texttt{pii-mcp}, \texttt{pii-web} and \texttt{pii-evolve}.
The daemon is not a fifth. \texttt{pii} embeds it in-process, so \texttt{pii daemon} runs the
control plane in the foreground and \texttt{pii -p}, \texttt{pii tui} and the command surface
either talk to a running daemon over the control socket or start an engine of their own. There
is one engine implementation and no client-side copy of the rules.

\begin{center}
\resizebox{\linewidth}{!}{%
\begin{tikzpicture}[node distance=8mm and 12mm]
	\node[b] (cli) {\texttt{pii-cli}};
	\node[b, right=of cli, xshift=6mm] (daemon) {daemon\\\scriptsize one writer};
	\node[s, right=of daemon, yshift=12mm] (store) {store.redb};
	\node[s, right=of daemon, yshift=-12mm] (dur) {durable/\\\scriptsize channel $\cdot$ timer};
	\node[b, below=of daemon] (sandbox) {sandbox\\\scriptsize landlock profile};
	\node[b, right=of sandbox] (worker) {worker\\\scriptsize agent argv};
	\node[b, left=of sandbox] (caps) {capability set\\\scriptsize ceiling $\cap$ node caps};
	\node[b, above=of daemon] (log) {daemon log\\\scriptsize JSON lines};

	\draw[f] (cli) -- node[p, above=2pt, midway]{socket} (daemon);
	\draw[f] (daemon) -- (store);
	\draw[f] (daemon) -- (dur);
	\draw[f] (daemon) -- (sandbox);
	\draw[f] (caps) -- (sandbox);
	\draw[f] (sandbox) -- (worker);
	\draw[f] (daemon) -- (log);
\end{tikzpicture}}
\captionof{figure}{One binary, one writer, a socket, and two independent bounds between the
engine and any agent: the process-wide sandbox and the per-node capability set.}
\end{center}

The two bounds are worth separating at the start because they answer different questions. The
capability set answers ``may this node do this kind of thing'' --- read files, write files,
reach the network, spawn a process. The sandbox answers ``may this process, at all'' --- and
it is enforced by the kernel, after the engine has decided. A node without \texttt{fs\_write}
cannot write whatever the agent believes about itself, and an agent that talks its way past
its own instructions still cannot write outside the granted paths.

\section{Boot order}

The order is the chapter's real subject. Each step below can refuse, and each refusal names
the thing that failed before anything has been bound, spawned or written.

\begin{enumerate}
	\item \textbf{Flags and config.} \texttt{pii daemon} takes \texttt{--store},
	\texttt{--config}, \texttt{--workspace}, \texttt{--sandbox}, \texttt{--log-file}. An
	unknown key in the config file is a load error, not a silent default: a typo'd grant must
	not become no grant.
	\item \textbf{The sandbox.} Under \texttt{--sandbox workspace}, confinement is applied and
	verified before the store is opened. Failure exits \texttt{4}.
	\item \textbf{The measurement switches.} \texttt{--observer} and \texttt{--audit} accept
	\texttt{on} or \texttt{off}; a typo is an error, never ``on''.
	\item \textbf{The log and the starting mark.} \texttt{daemon.starting} is written before
	anything that can fail, carrying the pid and the store parent, so a start that never serves
	is visible in the record rather than only in the supervisor's journal.
	\item \textbf{The durable claim.} \texttt{pii\_store::durable::claim} probes
	\texttt{channel.redb} and \texttt{timer.redb}; a store another process holds exits
	\texttt{3} with the store's path in the message.
	\item \textbf{The model.} No model configured exits \texttt{5}. The daemon will not start
	into a state where the first agent dispatch is the thing that discovers the omission.
	\item \textbf{The store and the engine.} One retention sweep runs here; a failure is
	reported as \texttt{store.retention.failed} and is not a refusal.
	\item \textbf{Channels, MQTT, cron, workers.} Each entry is validated; a cron entry whose
	graph does not start at a \texttt{timer} node is refused, because it would re-run as fast
	as it finishes.
	\item \textbf{The socket.} Bound last, after everything the daemon needs is known good.
\end{enumerate}

\begin{center}
\resizebox{\linewidth}{!}{%
\begin{tikzpicture}[node distance=8mm and 12mm]
	\node[b] (flags) {flags\\\scriptsize exit 2};
	\node[b, right=of flags] (sb) {sandbox\\\scriptsize exit 4};
	\node[b, right=of sb] (mark) {daemon.starting};
	\node[b, right=of mark] (claim) {durable claim\\\scriptsize exit 3};
	\node[b, below=of claim] (model) {model\\\scriptsize exit 5};
	\node[b, left=of model] (store) {store + engine\\\scriptsize sweep};
	\node[b, left=of store] (entries) {channels $\cdot$ cron\\\scriptsize validated};
	\node[b, left=of entries] (sock) {socket\\\scriptsize bound last};

	\draw[f] (flags) -- (sb);
	\draw[f] (sb) -- (mark);
	\draw[f] (mark) -- (claim);
	\draw[f] (claim) -- (model);
	\draw[f] (model) -- (store);
	\draw[f] (store) -- (entries);
	\draw[f] (entries) -- (sock);
	\draw[bk] (sb.north) -- ++(0,6mm) node[p, above=2pt, midway]{refuse};
\end{tikzpicture}}
\captionof{figure}{Every refusal happens before the socket exists, so a start that fails fails
where a person is watching a terminal rather than on the first request from a client.}
\end{center}

\begin{longtable}{@{}p{0.10\textwidth}p{0.34\textwidth}p{0.46\textwidth}@{}}
\toprule
\textbf{Exit} & \textbf{Condition} & \textbf{Message shape} \\
\midrule
\endhead
2 & bad flag or config & \texttt{config: <toml error>} or the usage line \\
3 & a durable store is held & \texttt{store <path> is locked: another daemon holds it} \\
4 & sandbox not enforceable & \texttt{sandbox: enforcement could not be verified} \\
5 & no model configured & \texttt{no model configured (set PII\_MODEL, or [core].model ...)} \\
\bottomrule
\end{longtable}

\begin{lstlisting}[language=bash]
pii daemon --store /var/lib/pii/store.redb \
           --workspace /projects/pi-improved \
           --sandbox workspace \
           --config ~/.pii/pii.toml \
           --log-file /var/log/pii/daemon.jsonl
\end{lstlisting}

The default store path is documented in the flags' own help as \texttt{<runtime
dir>/store.redb}, and the default workspace is the working directory. The defaults are the
thing to check when a daemon ``sees the wrong data'': the store path decides which history a
client is reading, and the workspace decides which files a sandboxed agent may write.

\section{One writer, enforced three times}

The store admits one process. That sentence is true at three levels, and each level catches a
different mistake.

\paragraph{The database file.} redb takes an exclusive lock on the file. A second process that
opens it either blocks or fails; either way two processes never interleave writes. The
single-writer rule in \texttt{pii-store} is about transactions inside one process; the file
lock is about processes.

\paragraph{The durable pair, claimed first.} The channel cursor and the timer wheel are
separate small databases under the state directory. Before binding anything, the daemon probes
both with a non-blocking lock, one attempt per store, no retry, no create. A store that is
free is left exactly as it was found; a store that is held is named. The probe deliberately
keeps no lock --- the owner opens the store immediately after --- so the claim is a check, not
a reservation.

\paragraph{The state directory.} \texttt{PII\_STATE\_DIR} is resolved before \texttt{HOME}, so
a second instance --- a test, a soak, an operator's second daemon --- opens its own store, its
own durable pair and its own audit log. Without that resolution order, a test started beside a
running daemon opens the operator's \texttt{channel.redb}, finds it held, and hangs before
binding its socket. The symptom is a test that never starts and a daemon log with nothing in
it.

\begin{center}
\resizebox{\linewidth}{!}{%
\begin{tikzpicture}[node distance=8mm and 12mm]
	\node[b] (a) {daemon A\\\scriptsize state dir S};
	\node[b, right=of a, yshift=10mm] (sa) {S/store.redb};
	\node[b, right=of a, yshift=-10mm] (da) {S/durable/};
	\node[b, below=of a] (b) {daemon B\\\scriptsize state dir S};
	\node[b, right=of b, yshift=-10mm] (ref) {refusal\\\scriptsize exit 3};
	\node[b, right=of b, yshift=10mm] (sb) {S/store.redb};

	\draw[f] (a) -- (sa);
	\draw[f] (a) -- (da);
	\draw[bk] (b) -- node[p, below=2pt, midway]{claim} (ref);
	\draw[sf, dashed] (sb) -- (sa);
	\node[p, right=of sa, xshift=2mm, font=\scriptsize\itshape, align=left]
{held by A};
\end{tikzpicture}}
\captionof{figure}{The second daemon is refused by name rather than queued behind a lock,
because waiting produced a hang with no message that reads as a broken binary.}
\end{center}

The rule has a cost the operator must know: two daemons cannot share a state directory, so two
concurrent soaks or two test suites need two directories. The failure when they do not is loud
now, which is the point; before it was loud, it was a flaky test attributed to the code under
test.

\section{Two directories that mean different things}

The runtime directory holds things that are allowed to vanish on reboot:
\texttt{control.sock}, the graphs a daemon loads, \texttt{crashes/}. On this host it is a
per-user directory under \texttt{/run/user}. The state directory holds things that must
survive: the run store, the durable pair, the audit log.

Keeping them apart is a decision with a visible consequence. A store under the runtime
directory would be emptied by a reboot, and a reboot would therefore silently delete every run
history and every pending timer. The durable directory is resolved as
\texttt{[core].durable\_dir} when set, else \texttt{<state\_dir>/durable}, and never defaults
to the runtime directory. The comment in \texttt{daemon/mod.rs} says why in one clause: the
runtime dir is a tmpfs.

\section{The sandbox is the real bound}

\texttt{crates/pii-tools/src/sandbox.rs} implements two profiles and one rule about them.

\begin{longtable}{@{}p{0.14\textwidth}p{0.78\textwidth}@{}}
\toprule
\textbf{Profile} & \textbf{What is in force} \\
\midrule
\endhead
\texttt{workspace} & WRITES confined to the workspace, the store directory, the runtime directory,
\texttt{\$TMPDIR}, the container runtime state dirs and \texttt{/dev/shm}, plus the device
files \texttt{/dev/null}, \texttt{/dev/zero}, \texttt{/dev/urandom}, \texttt{/dev/tty}. READS
are not confined: read-only access to \texttt{/} is granted, so \texttt{/etc}, \texttt{HOME}
and every other
readable path stay readable. \\
\texttt{off} & no filesystem confinement in force; the explicit opt-out. Reported loudly at start,
in the log, in the audit trail and by \texttt{pii health}. \\
\bottomrule
\end{longtable}

The asymmetry --- writes confined, reads not --- is stated in the same words in four places:
the profile's grant text, the startup event, \texttt{pii health}, and the deploy runbook. One
text, so a reader cannot get a more generous impression from one surface than another. The
test suite pins the wording, which is unusual and deliberate: a permission model that is
described differently in two places is two permission models.

Applying the profile is \texttt{apply\_workspace}, and it is fail-closed in a way worth
reading line by line.

\begin{enumerate}
	\item The ruleset handles every filesystem access the ABI knows, and is created before any
	rule is added. A failure here is an error, not a partial ruleset.
	\item Every read-only path (\texttt{/} and \texttt{/dev}) is added as read access. A
	missing path is skipped rather than fatal, because a container may not have it.
	\item Every writable path is added with full directory rights --- except device nodes,
	which get file-content rights only. Granting directory rights on a device file makes the
	kernel downgrade the whole ruleset to partially enforced, which would silently weaken
	everything else.
	\item The ruleset is applied to the process. If the resulting status is not fully enforced,
	or \texttt{no\_new\_privs} is not set, the function returns an error. The caller exits
	\texttt{4}.
	\item An \emph{out-of-grant} probe writes a file in a directory that must not be writable.
	If the write succeeds, confinement is not real and the function fails.
	\item An \emph{in-grant} probe writes a file in a directory that must be writable. If that
	write fails, the grant is inert --- every rule was skipped or downgraded --- and the
	function fails.
\end{enumerate}

\begin{center}
\resizebox{\linewidth}{!}{%
\begin{tikzpicture}[node distance=8mm and 12mm]
	\node[b] (rules) {ruleset\\\scriptsize ro: / and /dev};
	\node[b, right=of rules] (rw) {rw: listed dirs\\\scriptsize files get content rights};
	\node[b, below=of rw] (restrict) {\texttt{restrict\_self}};
	\node[v, below=of restrict] (status) {enforced?};
	\node[b, left=of status] (exit) {exit 4};
	\node[b, below=of status] (out) {out-grant probe\\\scriptsize must fail};
	\node[b, below=of out] (in) {in-grant probe\\\scriptsize must succeed};
	\node[b, left=of in] (serve) {serve};

	\draw[f] (rules) -- (rw);
	\draw[f] (rw) -- (restrict);
	\draw[f] (restrict) -- (status);
	\draw[f] (status) -- node[p, above=2pt, midway]{yes} (out);
	\draw[bk] (status) -- node[p, above=2pt, midway]{no} (exit);
	\draw[f] (out) -- (in);
	\draw[f] (in) -- node[p, above=2pt, midway]{ok} (serve);
	\draw[bk] (in.south) |- node[p, below=2pt, midway]{inert} (exit.south);
\end{tikzpicture}}
\captionof{figure}{Six checks, two of them probes in opposite directions, because without the
in-grant probe a ruleset whose every rule was skipped would pass the out-grant test while
granting nothing, and every later refusal becomes a coincidence.}
\end{center}

The two probes are the part that survives review. A single out-grant probe proves that
something is denied; it cannot distinguish ``confined correctly'' from ``confined to
nothing''. The in-grant probe is the control that makes the first probe meaningful, and it is
the reason a report of a green sandbox on a new host can be trusted.

\paragraph{Two grants that are not obvious.} Podman keeps its graph root under
\texttt{HOME/.local/share/containers} and its runroot and locks under
\texttt{XDG\_RUNTIME\_DIR/containers} and \texttt{XDG\_RUNTIME\_DIR/libpod}. A sandboxed
daemon without those could not inspect a present image: the measured symptom was
\texttt{EACCES} on \texttt{bolt\_state.db} and \texttt{alive.lck}. The fourth entry is
\texttt{/dev/shm}, because podman's rootless lock file lives there and the profile grants
\texttt{/dev} read-only. The podman \emph{socket} is deliberately not a grant: a process
reaches it by \texttt{connect()}, which landlock does not mediate, so adding it would change
the grant text without changing what is allowed. A test asserts the text does not mention it,
which keeps the documentation and the mechanism from drifting apart.

\paragraph{The forced-failure switch.} \texttt{PII\_SANDBOX\_FORCE\_UNSUPPORTED=1} makes
\texttt{apply\_workspace} return the unverifiable error immediately. It exists so the exit-4
path can be exercised on a host where landlock \emph{does} work --- a refusal path that is
never run is a refusal path that does not work.

\section{The capability ceiling}

The sandbox is process-wide and coarse. The capability set is per-run and per-node, and it is
the bound an agent's own instructions cannot argue with.

A capability set is twelve bits plus five scopes. The bits are stable forever, because a
worker or a plugin references them without depending on the engine:

\begin{longtable}{@{}p{0.16\textwidth}p{0.06\textwidth}p{0.66\textwidth}@{}}
\toprule
\textbf{Kind} & \textbf{Bit} & \textbf{What it permits} \\
\midrule
\endhead
\texttt{fs\_read} & 1 & read files within an allow path scope \\
\texttt{fs\_write} & 2 & write files within an allow path scope \\
\texttt{net} & 4 & reach the network \\
\texttt{proc} & 8 & spawn or talk to processes \\
\texttt{hw} & 16 & touch hardware \\
\texttt{secrets} & 32 & access secrets \\
\texttt{memory\_write} & 64 & write into the agent's memory \\
\texttt{spawn} & 128 & spawn subprocesses \\
\texttt{memory\_read} & 256 & read from the agent's memory \\
\texttt{post} & 512 & post into a configured channel \\
\texttt{mqtt\_read} & 1024 & list and subscribe to a house's devices \\
\texttt{mqtt\_write} & 2048 & command a device through its allowlist \\
\bottomrule
\end{longtable}

The only operation on two sets is intersection. There is deliberately no union: a capability
set can be narrowed at every hop and widened at none. \texttt{CapSet::intersect} ANDs the bits
and intersects every scope; \texttt{narrow\_kinds} lowers the bits while keeping the inherited
path and host scopes, which is what a node that declares fewer capabilities than its run gets.

On top of the per-run set sits a ceiling from configuration: \texttt{[grant\_ceiling]} names
the capability kinds no graph or node may request beyond. An empty list means no ceiling; a
name the engine does not know is an invalid-config error rather than a silently ignored entry.
The check runs where the node is built, so a graph asking for more than the ceiling is refused
before it runs, and the refusal names the capability.

Path scopes are deny-wins. \texttt{PathScope::permits} checks the deny list first, so a path
inside an allow prefix but inside a deny prefix is refused. The heal path uses that to keep a
patch away from \texttt{specs}, \texttt{xtask}, \texttt{quality-gate.sh}, \texttt{graphs},
\texttt{clippy.toml}, \texttt{rust-toolchain.toml} and \texttt{.github}: a heal worktree is a
full checkout, and without the deny list a patch could delete the assertions that prove it
worked.

\begin{lstlisting}[language=bash]
[core]
model = "llamacpp-server:granite-4.2-3b"

# the ceiling: no graph or node may request beyond these kinds
grant_ceiling = ["fs_read", "fs_write", "proc", "spawn", "net"]

[net]
grants = ["10.10.10.217:8082"]
allow_private = true

[[model_tier]]
name = "small"
model = "llamacpp-server:granite-4.2-3b"
tools = ["read_file", "grep"]
max_tokens = 2048
\end{lstlisting}

\paragraph{Why the two bounds are both needed.} A capability set with no sandbox is advice to
a process that has already been spawned by an engine with the operator's rights. A sandbox
with no capability set confines a process but gives every node the same power inside the
fence. The pairing is what makes ``this node may write, that node may not, and neither may
write outside the workspace'' expressible.

\section{The socket, and why clients are thin}

Every client reaches the daemon through one unix socket, \texttt{control.sock}, under the
runtime directory. The socket is the boundary that keeps a second copy of the rules from
existing: \texttt{pii run ls}, \texttt{pii health}, \texttt{pii queue add} and \texttt{pii
skill show} all become a request with a method and a payload, and the engine answers. A client
cannot be more permissive than the daemon, because it has no engine of its own to be
permissive with.

The socket's location is computed once, from the runtime directory, and \texttt{pii-uds} owns
that computation --- one default socket, so a CLI started from a plain shell with no
environment finds the daemon that an operator started from a session with one. Before that,
\texttt{pii health} answered for a daemon that did not exist or failed to find one that did,
and the difference was an unset variable nobody could see.

This is also why a long-running command does not hold a store lock. The client sends a request
and waits; the work happens in the daemon, in the daemon's process, under the daemon's sandbox
and capability ceiling. The \texttt{--store} flag on a one-shot command exists for the case
where no daemon is running, and then the process really does open the file, which is why a
second one-shot beside a running daemon can collide where a socket request would not.

\section{What the daemon runs}

Beyond runs, the daemon owns a small set of long-lived things, each with its own failure mode:

\begin{itemize}
	\item \textbf{Channels} --- one runtime per configured channel, with a poll interval, a
	poll backoff and a graph to start per inbound message. The default poll interval is 1\,500
	ms because Discord answers a message read at roughly five requests per five seconds per
	route; at 500 ms the daemon collected 173 \texttt{channel.poll.failed 429} events in a few
	hours and never read a message.
	\item \textbf{MQTT connections} --- one broker connection each, subscribing to topic
	filters and starting a graph per message, with a bounded inbound queue that drops and
	counts beyond its cap. A name used by both a channel and an MQTT connection is a config
	error, not a last-writer-wins.
	\item \textbf{Cron entries} --- display-only schedules; the actual firing comes from the
	graph's entry \texttt{timer} node. The entry kind is checked at start, because a non-timer
	entry would re-run the instant it finishes.
	\item \textbf{Workers} --- declared in config as a program, its argv, extra environment and
	an in-flight cap, so a bounded number of external processes exist at once.
	\item \textbf{Timers} --- the durable wheel, rearmed from the store at boot.
	\item \textbf{Claim recovery} --- \texttt{requeue\_dead\_claims} at boot returns rows whose
	claimant is gone to \texttt{Pending}, so a killed daemon does not leave work claimed
	forever.
\end{itemize}

\section{Refusals that name the thing}

A refusal is only useful if it names what failed and where. The daemon's refusals follow one
shape: the object, the reason, and the identifier.

\begin{lstlisting}[language=bash]
store /root/.local/state/pii/durable/channel.redb is locked: another daemon holds it
cron 'nightly': graph 'night' must start at a timer node (entry is gate); it would re-run
  as fast as it finishes
config: [core].ollama_host is retired (spec 129) - set `[llama] host` for the
  llama.cpp server instead
config: duplicate [[model_tier]] name 'small'
sandbox: enforcement could not be verified
\end{lstlisting}

Two of those deserve comment. The retired \texttt{ollama\_host} key is refused \emph{by name}
rather than as an unknown field, because the generic message would send an operator looking
for a typo. And the duplicate tier name is refused at parse time, so a node asking for the
tier gets one provider rather than whichever entry happened to win.

\section{What an operator should check first}

When a daemon is suspected rather than proven wrong, four checks answer most questions without
reading code:

\begin{enumerate}
	\item \texttt{pii health --json}: the state of each subsystem, the store's path and
	per-table counts, and the mounts the process actually has.
	\item The log's first line: \texttt{daemon.starting} carries the pid and the store parent,
	so the process and the directory it opened are in the record.
	\item The sandbox line in \texttt{health}: \texttt{sandbox=workspace (landlock)} with the
	grants spelled out, or the loud \texttt{off} text.
	\item The store's parent directory, compared with \texttt{PII\_STATE\_DIR} and
	\texttt{HOME}: a daemon reading the wrong store looks exactly like a daemon with amnesia.
\end{enumerate}

\section{Failure modes worth recognising}

\begin{enumerate}
	\item \textbf{A test daemon that hangs on start.} Symptom: no socket, no log line after
	\texttt{daemon.starting}, a test that times out. Cause: it inherited the operator's
	\texttt{PII\_STATE\_DIR} and is waiting on a held \texttt{channel.redb}. Fix: give the test
	its own state directory.
	\item \textbf{A gate that fails on a channel lock.} Symptom: \texttt{channel.redb is
	locked: another daemon holds it} inside a gate run. Cause: the gate inherited a state dir
	from the shell that started it. The gate now removes \texttt{PII\_STATE\_DIR} and
	\texttt{PII\_RUNTIME\_DIR} from the child environment.
	\item \textbf{A sandbox that refuses every start.} Symptom: exit \texttt{4} on a host where
	landlock should work. Read the two probe outcomes: an out-grant write that succeeded means
	the ruleset did not apply; an in-grant write that failed means the grant list is wrong,
	usually a missing directory the process cannot create.
	\item \textbf{A podman run that cannot see its own image.} Symptom: \texttt{EACCES} on
	\texttt{bolt\_state.db}. Cause: the container state dirs are not in the grant list.
	\item \textbf{A node that ``has no write tool''.} Symptom: the agent stops and reports that
	no write or execute tool is available. Cause: the node declared too few capabilities --- a
	coder without \texttt{fs\_write} cannot implement anything. The fix is in the graph, not in
	the sandbox.
	\item \textbf{A cron entry that spins.} Symptom: a graph that starts again immediately
	after it finishes. Cause: the entry node is not a \texttt{timer}. The daemon refuses that
	at start now.
\end{enumerate}
